\subsubsection{Threshold difference and composition of shells}
\label{mg07:umbral}

The enlargement of a discrete cell preserves every stratum appearing
when its size increases by one unit. The threshold difference expresses
this operation on the cardinality reader and tracks its behaviour under
addition and multiplication before a metric realization is introduced.

\begin{definition}[Threshold difference]
Let \(R\) be a commutative ring with identity and let
\(f:\mathbb Z_{\geq0}\to R\). Define
\begin{equation}
 (D_\#f)(n)=f(n+1)-f(n).
 \label{mg07:definicion}
\end{equation}
For an integer reader of nested finite cells, this difference records
the cardinality of the shell added in the step \(n\to n+1\).
\end{definition}

\begin{proposition}[Additivity and exact product law]
For \(f,g:\mathbb Z_{\geq0}\to R\),
\begin{align}
 D_\#(f+g)(n)&=D_\#f(n)+D_\#g(n),
 \label{mg07:aditividad}\\
 D_\#(fg)(n)&=f(n+1)D_\#g(n)+g(n)D_\#f(n)
 \label{mg07:producto-desplazado}\\
 &=f(n)D_\#g(n)+g(n)D_\#f(n)
   +D_\#f(n)D_\#g(n).
 \label{mg07:producto}
\end{align}
\end{proposition}

\begin{proof}
The difference of a sum separates term by term. For the product,
adding and subtracting \(f(n+1)g(n)\) gives
\[
 \begin{split}
 f(n+1)g(n+1)-f(n)g(n)
 &=f(n+1)\bigl(g(n+1)-g(n)\bigr)\\
 &\quad+g(n)\bigl(f(n+1)-f(n)\bigr).
 \end{split}
\]
This proves \eqref{mg07:producto-desplazado}. Substitution of
\(f(n+1)=f(n)+D_\#f(n)\) yields \eqref{mg07:producto}, including
its complete bilinear term.
\end{proof}

If \(A_n\subseteq A_{n+1}\) and \(B_n\subseteq B_{n+1}\) are
finite sets, the shell of the product decomposes into three disjoint
parts:
\[
 \begin{split}
 (A_{n+1}\times B_{n+1})\setminus(A_n\times B_n)
 &=\bigl(A_n\times(B_{n+1}\setminus B_n)\bigr)\\
 &\quad\sqcup\bigl((A_{n+1}\setminus A_n)\times B_n\bigr)\\
 &\quad\sqcup\bigl((A_{n+1}\setminus A_n)
                 \times(B_{n+1}\setminus B_n)\bigr).
 \end{split}
\]
Taking cardinalities gives exactly \eqref{mg07:producto}. The
bilinear term counts positions new in both factors; it resides in
the intersection of the two enlargements, with incidence preserved.

\begin{proposition}[Binomial law for the dimensional shell]
For \(d\geq1\), let \(f_d(n)=n^d\). Then
\begin{equation}
 D_\#f_d(n)=(n+1)^d-n^d
 =\sum_{k=0}^{d-1}\binom dk n^k
 =\sum_{r=1}^{d}\binom dr n^{d-r}.
 \label{mg07:binomio}
\end{equation}
\end{proposition}

\begin{proof}
The binomial expansion of \((n+1)^d\) contains \(n^d\) and the
remaining \(d\) terms in the first sum. Substituting \(r=d-k\)
gives the second sum. In the support \(\{1,\ldots,n+1\}^d\),
each added position has a nonempty set of coordinates equal to
\(n+1\). If that set has \(r\) elements, there are \(\binom dr\)
choices of positions and \(n^{d-r}\) choices of previous coordinate
values. These strata are disjoint and cover the shell, also proving
the formula by finite incidence.
\end{proof}

For \(d=3\), the three strata have cardinalities \(3n^2\),
\(3n\) and \(1\). At the nonadic--decimal step, \(n=9\), the
operation gives \(243+27+1=271\); excluding the final vertex leaves
\(270\) positions. The following cubic completion makes explicit
the union of this shell with the \(729\) previous positions. The
threshold law thus preserves both the total and its distribution by
the number of new coordinates. Subsequent lattice realizations
transport their incidences through the maps specific to each
construction.
